\subsection{Nguyên lý Infrastructure as Code và Quản lý gói ứng dụng}

\subsubsection{Nguyên lý Infrastructure as Code (IaC)} Trong vận hành hệ thống truyền thống, việc thiết lập hạ tầng thường được thực hiện thủ công hoặc thông qua các tập lệnh mang tính mệnh lệnh. Phương pháp này tồn tại nhiều rủi ro về tính không đồng nhất giữa các môi trường và khó khăn trong việc mở rộng. Để khắc phục, đồ án áp dụng nguyên lý Infrastructure as Code (IaC) – một phương pháp tiếp cận trong đó việc quản lý và cung cấp tài nguyên máy tính được thực hiện thông qua các tệp định nghĩa máy có thể đọc được.

IaC hoạt động dựa trên mô hình Khai báo. Thay vì chỉ định từng bước "làm thế nào" để đạt được trạng thái mong muốn, kỹ sư chỉ cần định nghĩa "trạng thái cuối cùng" của hệ thống. Nền tảng điều phối sẽ tự động tính toán và thực thi các thay đổi cần thiết để đưa hệ thống thực tế về trạng thái đã định nghĩa. Cách tiếp cận này biến hạ tầng thành phần mềm, cho phép áp dụng các quy trình kiểm soát phiên bản, kiểm thử và tích hợp liên tục vào việc quản lý hạ tầng.

\subsubsection{Cơ chế Quản lý gói ứng dụng} Mặc dù IaC giúp định nghĩa hạ tầng dưới dạng mã, nhưng trong các hệ thống phân tán phức tạp, số lượng các tệp cấu hình (như các file YAML trong Kubernetes) có thể tăng lên rất nhanh và trở nên rời rạc, khó kiểm soát. Để giải quyết vấn đề tổ chức và phân phối các tệp cấu hình này, khái niệm Quản lý gói cho hạ tầng được áp dụng.

Về bản chất, quản lý gói trong ngữ cảnh hạ tầng đám mây giải quyết ba bài toán cốt lõi: 
\begin{itemize} 
    \item \textbf{Đóng gói và Trừu tượng hóa:} Thay vì quản lý hàng loạt các tài nguyên riêng lẻ (như mạng, lưu trữ, tính toán), cơ chế này cho phép gom nhóm tất cả các thành phần liên quan của một ứng dụng thành một đơn vị duy nhất (một "gói"). Điều này giúp người vận hành nhìn nhận hệ thống ở mức độ dịch vụ tổng thể thay vì các thành phần vụn vặt. 
    \item \textbf{Cơ chế mẫu:} Quản lý gói tách biệt giữa "cấu trúc hạ tầng" và "dữ liệu cấu hình". Thông qua việc sử dụng các biểu mẫu (templates) và biến số (variables), cùng một gói cấu hình có thể được tái sử dụng linh hoạt cho nhiều môi trường khác nhau (Development, Staging, Production) mà không cần sao chép hay chỉnh sửa mã nguồn gốc. 
    \item \textbf{Quản lý phiên bản và vòng đời:} Các gói cấu hình được đánh số phiên bản cụ thể, cho phép hệ thống dễ dàng thực hiện việc nâng cấp hoặc quay lui toàn bộ ứng dụng về trạng thái ổn định trước đó chỉ bằng một thao tác, đảm bảo tính toàn vẹn của hệ thống. 
\end{itemize}